\textbf{Phần 1. Trắc nghiệm lựa chọn}
\setcounterpageref{ex}{0}

\begin{ex} %[5D7N44]
Tỉ số của hai số $5$ và $9$ là
\choice
{\True $\dfrac{5}{9}$}
{$\dfrac{9}{5}$}
{$5 \times 9$}
{$45$}
\loigiai{
Tỉ số của hai số $5$ và $9$ là $\dfrac{5}{9}$.\\
\textbf{Đáp án của câu hỏi là A}
}
\end{ex}

\begin{ex} %[5D7N44]
Tỉ số phần trăm $48\%$ được viết dưới dạng số thập phân là
\choice
{$4{,}8$}
{$0{,}048$}
{\True $0{,}48$}
{$48$}
\loigiai{
Ta có: $48\% = \dfrac{48}{100} = 0{,}48$.\\
\textbf{Đáp án của câu hỏi là C}
}
\end{ex}

\begin{ex} %[5D7H44]
Trên bản đồ tỉ lệ $1 : 50\,000$, khoảng cách giữa hai địa điểm đo được là $4\text{ cm}$. Khoảng cách thật ngoài thực tế là
\choice
{$200\text{ m}$}
{$20\text{ km}$}
{$200\text{ km}$}
{\True $2\text{ km}$}
\loigiai{
Khoảng cách thực tế: $4 \times 50\,000 = 200\,000\text{ cm} = 2\text{ km}$.\\
\textbf{Đáp án của câu hỏi là D}
}
\end{ex}

\begin{ex} %[5D7H44]
Tổng của hai số là $80$, tỉ số của hai số là $\dfrac{2}{3}$. Số lớn là
\choice
{$32$}
{\True $48$}
{$50$}
{$60$}
\loigiai{
Tổng số phần: $2 + 3 = 5$ (phần).\\
Số lớn là: $80 : 5 \times 3 = 48$.\\
\textbf{Đáp án của câu hỏi là B}
}
\end{ex}

\begin{ex} %[5D7H44]
Hiệu của hai số là $42$, tỉ số của hai số là $\dfrac{3}{5}$. Số bé là
\choice
{$21$}
{$105$}
{\True $63$}
{$84$}
\loigiai{
Hiệu số phần: $5 - 3 = 2$ (phần).\\
Số bé là: $42 : 2 \times 3 = 63$.\\
\textbf{Đáp án của câu hỏi là C}
}
\end{ex}

\begin{ex} %[5D7H44]
Một cửa hàng bán được $450\text{ kg}$ gạo tẻ và $150\text{ kg}$ gạo nếp. Tỉ số phần trăm của số gạo nếp so với tổng số gạo đã bán là
\choice
{$30\%$}
{\True $25\%$}
{$33{,}3\%$}
{$20\%$}
\loigiai{
Tổng số gạo: $450 + 150 = 600\text{ (kg)}$.\\
Tỉ số phần trăm gạo nếp: $150 : 600 = 0{,}25 = 25\%$.\\
\textbf{Đáp án của câu hỏi là B}
}
\end{ex}

\begin{ex} %[5D7V44]
Một mảnh đất hình chữ nhật có chu vi $160\text{ m}$, chiều rộng bằng $\dfrac{3}{5}$ chiều dài. Diện tích mảnh đất đó là
\choice
{\True $1\,500\text{ m}^2$}
{$1\,200\text{ m}^2$}
{$1\,600\text{ m}^2$}
{$6\,000\text{ m}^2$}
\loigiai{
Nửa chu vi: $160 : 2 = 80\text{ (m)}$.\\
Chiều dài: $80 : (3 + 5) \times 5 = 50\text{ (m)}$.\\
Chiều rộng: $80 - 50 = 30\text{ (m)}$.\\
Diện tích: $50 \times 30 = 1\,500\text{ (m}^2\text{)}$.\\
\textbf{Đáp án của câu hỏi là A}
}
\end{ex}

\begin{ex} %[5D7V44]
Một người gửi tiết kiệm $120\,000\,000$ đồng với lãi suất $0{,}5\%$ một tháng. Sau một tháng cả tiền gửi và tiền lãi là
\choice
{$120\,500\,000$ đồng}
{$126\,000\,000$ đồng}
{$125\,000\,000$ đồng}
{\True $120\,600\,000$ đồng}
\loigiai{
Tiền lãi 1 tháng: $120\,000\,000 \times 0{,}5\% = 600\,000$ (đồng).\\
Tổng số tiền: $120\,000\,000 + 600\,000 = 120\,600\,000$ (đồng).\\
\textbf{Đáp án của câu hỏi là D}
}
\end{ex}

\begin{ex} %[5D7V44]
Năm nay tuổi bố gấp $4$ lần tuổi con. Biết $4$ năm trước bố hơn con $30$ tuổi. Tuổi bố hiện nay là
\choice
{$36$ tuổi}
{$44$ tuổi}
{\True $40$ tuổi}
{$48$ tuổi}
\loigiai{
Hiệu số tuổi không đổi, hiện nay bố vẫn hơn con $30$ tuổi.\\
Hiệu số phần: $4 - 1 = 3$ (phần).\\
Tuổi bố hiện nay là: $30 : 3 \times 4 = 40$ (tuổi).\\
\textbf{Đáp án của câu hỏi là C}
}
\end{ex}

\begin{ex} %[5D7V44]
Một người mua một chiếc xe đạp giá $2\,000\,000$ đồng. Sau đó bán lại bị lỗ $5\%$. Người đó đã bán chiếc xe với giá là
\choice
{\True $1\,900\,000$ đồng}
{$1\,850\,000$ đồng}
{$1\,950\,000$ đồng}
{$1\,800\,000$ đồng}
\loigiai{
Số tiền lỗ: $2\,000\,000 \times 5\% = 100\,000$ (đồng).\\
Giá bán: $2\,000\,000 - 100\,000 = 1\,900\,000$ (đồng).\\
\textbf{Đáp án của câu hỏi là A}
}
\end{ex}

\begin{ex} %[5D7V44]
Trên bản đồ tỉ lệ $1 : 2\,500$, một mảnh ruộng hình vuông có chu vi là $16\text{ cm}$. Diện tích thực tế của mảnh ruộng đó là
\choice
{$1\text{ ha}$}
{\True $1\text{ ha}$}
{$10\text{ ha}$}
{$0{,}1\text{ ha}$}
\loigiai{
Cạnh trên bản đồ: $16 : 4 = 4\text{ (cm)}$.\\
Cạnh thực tế: $4 \times 2\,500 = 10\,000\text{ cm} = 100\text{ m}$.\\
Diện tích thực tế: $100 \times 100 = 10\,000\text{ m}^2 = 1\text{ ha}$.\\
\textbf{Đáp án của câu hỏi là B}
}
\end{ex}

\begin{ex} %[5D7C44]
Một cửa hàng nhân ngày Quốc tế Thiếu nhi đã hạ giá các mặt hàng $15\%$, tuy vậy cửa hàng vẫn còn lãi $10{,}5\%$ so với giá vốn. Hỏi nếu không hạ giá thì cửa hàng lãi bao nhiêu phần trăm so với giá vốn?
\choice
{$25\%$}
{$28\%$}
{$32\%$}
{\True $30\%$}
\loigiai{
Coi giá vốn là $100\%$.\\
Giá bán ngày lễ là: $100\% + 10{,}5\% = 110{,}5\%$ (giá vốn).\\
Giá bán ngày lễ bằng $100\% - 15\% = 85\%$ giá định bán.\\
Giá định bán so với giá vốn: $110{,}5\% : 85 \times 100 = 130\%$.\\
Nếu không hạ giá thì lãi: $130\% - 100\% = 30\%$.\\
\textbf{Đáp án của câu hỏi là D}
}
\end{ex}

\textbf{Phần 2. Trắc nghiệm đúng sai}
\setcounterpageref{ex}{0}

\begin{ex} %[5D7H44]
Một nông trại nuôi bò và dê với tổng cộng $450$ con, số bò bằng $\dfrac{4}{5}$ số dê. Xét tính đúng sai:
\choiceTF
{\True Số bò chiếm $4$ phần, số dê chiếm $5$ phần bằng nhau.}
{Số bò là $250$ con.}
{\True Số dê là $250$ con.}
{\True Số bò chiếm khoảng $44{,}4\%$ tổng đàn gia súc.}
\loigiai{
\begin{itemize}
    \item Tỉ lệ bò và dê: $4 : 5$. Đúng.
    \item Số bò: $450 : (4+5) \times 4 = 200$ con $\neq 250$ con. Sai.
    \item Số dê: $450 - 200 = 250$ con. Đúng.
    \item Tỉ số phần trăm bò: $200 : 450 \approx 44{,}44\%$. Đúng.
\end{itemize}
\begin{itemize}
    \item a) \textbf{Đúng}.
    \item b) \textbf{Sai}.
    \item c) \textbf{Đúng}.
    \item d) \textbf{Đúng}.
\end{itemize}
}
\end{ex}

\begin{ex} %[5D7V44]
Một mảnh đất hình chữ nhật có chu vi $120\text{ m}$, chiều rộng bằng $\dfrac{2}{3}$ chiều dài. Người ta dùng $25\%$ diện tích để trồng cây ăn quả, $35\%$ diện tích để làm ao cá, phần còn lại trồng rau. Xét tính đúng sai:
\choiceTF
{\True Diện tích mảnh đất là $864\text{ m}^2$.}
{\True Diện tích trồng cây ăn quả là $216\text{ m}^2$.}
{Diện tích làm ao cá là $310\text{ m}^2$.}
{\True Diện tích trồng rau là $345{,}6\text{ m}^2$.}
\loigiai{
\begin{itemize}
    \item Nửa chu vi: $60\text{ m}$. Chiều dài: $36\text{ m}$, chiều rộng: $24\text{ m}$. Diện tích: $36 \times 24 = 864\text{ m}^2$. Đúng.
    \item Trồng cây: $864 \times 25\% = 216\text{ m}^2$. Đúng.
    \item Ao cá: $864 \times 35\% = 302{,}4\text{ m}^2 \neq 310\text{ m}^2$. Sai.
    \item Trồng rau: $864 \times 40\% = 345{,}6\text{ m}^2$. Đúng.
\end{itemize}
\begin{itemize}
    \item a) \textbf{Đúng}.
    \item b) \textbf{Đúng}.
    \item c) \textbf{Sai}.
    \item d) \textbf{Đúng}.
\end{itemize}
}
\end{ex}

\textbf{Phần 3. Trả lời ngắn}
\setcounterpageref{ex}{0}

\begin{ex} %[5D7H44]
Tìm hai số biết tổng là $198$ và tỉ số là $\dfrac{4}{5}$. Tìm số bé.
\shortans{88}
\loigiai{
\textbf{Đáp số:} $88$\\
$198 : (4 + 5) \times 4 = 88$.
}
\end{ex}

\begin{ex} %[5D7H44]
Hiệu của hai số là $64$, số lớn gấp $5$ lần số bé. Tìm số lớn.
\shortans{80}
\loigiai{
\textbf{Đáp số:} $80$\\
$64 : (5 - 1) \times 5 = 80$.
}
\end{ex}

\begin{ex} %[5D7V44]
Tìm $35\%$ của $640\text{ kg}$.
\shortans{224}
\loigiai{
\textbf{Đáp số:} $224$\\
$640 \times 35 : 100 = 224\text{ kg}$.
}
\end{ex}

\begin{ex} %[5D7V44]
Một trường tiểu học có $600$ học sinh, trong đó có $318$ học sinh nữ. Tính tỉ số phần trăm của số học sinh nữ so với toàn trường. (Chỉ ghi số).
\shortans{53}
\loigiai{
\textbf{Đáp số:} $53$\\
$318 : 600 = 53\%$.
}
\end{ex}

\begin{ex} %[5D7V44]
Trên bản đồ tỉ lệ $1 : 20\,000$, một con kênh dài $15\text{ cm}$. Tính độ dài thật của con kênh đó theo ki-lô-mét.
\shortans{3}
\loigiai{
\textbf{Đáp số:} $3$\\
$15 \times 20\,000 = 300\,000\text{ cm} = 3\text{ km}$.
}
\end{ex}

\begin{ex} %[5D7C44]
Một cửa hàng bán được $480\text{ kg}$ gạo tẻ và nếp, thu được $9\,600\,000$ đồng. Biết giá $1\text{ kg}$ gạo nếp đắt hơn gạo tẻ là $5\,000$ đồng và số gạo tẻ gấp $3$ lần số gạo nếp. Tính giá tiền $1\text{ kg}$ gạo tẻ theo đơn vị nghìn đồng.
\shortans{18,75}
\loigiai{
\textbf{Đáp số:} $18{,}75$\\
Số gạo nếp: $480 : (3 + 1) = 120\text{ kg}$.\\
Số gạo tẻ: $480 - 120 = 360\text{ kg}$.\\
Nếu 120 kg nếp bán bằng giá tẻ thì tổng số tiền giảm đi: $120 \times 5\,000 = 600\,000$ đồng.\\
Tổng tiền khi 480 kg đều là gạo tẻ: $9\,600\,000 - 600\,000 = 9\,000\,000$ đồng.\\
Giá 1 kg gạo tẻ: $9\,000\,000 : 480 = 18\,750$ đồng $= 18{,}75$ nghìn đồng.
}
\end{ex}
